\subsection{局部多面体集的 Stokes 公式}

本号中，X 表示类 \(C^{r}\)（\(r \geqslant 2\)）的有限维纯实流形；我们假定 X 是分离的。\footnote{读者若对更一般的情形感兴趣，可参阅 H. Whitney, \emph{Geometric Integration Theory}, 第 III 章 § 18（Princeton Univ. Press, 1957）。}

11.3.1. 设 A 是有限维实向量空间的一个子集。如果 A 是有限个闭半空间的有限交的并，则称 A 是多面体的。

称 X 的子集 A 是局部多面体的，如果对于每个 \(x \in X\)，都存在 X 在 x 处的图 \(c = (\mathrm{U}, \varphi, \mathrm{E})\) 以及 E 的一个多面体子集 \(A_{c}\)，使得 \(\varphi(\mathrm{A} \cap \mathrm{U}) = \varphi(\mathrm{U}) \cap \mathrm{A}_{c}\)。X 的一个片段是局部多面体子集。

11.3.2. 设 A 是 X 的闭子集，令 \(\mathrm{Fr}(\mathrm{A}) = \mathrm{A} - \mathring{\mathrm{A}}\) 为其边界。称 \(x \in \mathrm{Fr}(\mathrm{A})\) 为正则点，如果存在 x 的一个开邻域 U，使得 \(A \cap U\) 是流形 U 的一个片段（此时 x 属于 \(A \cap U\) 的边界）。我们把 \(\partial A\) 记作 \(\mathrm{Fr}(\mathrm{A})\) 的正则点集合，并称其为 A 的正则边界（或简称边界）。集合 \(A' = \mathring{A} \cup \partial A\) 是 X 的一个片段，其边界为 \(\partial A\)。

11.3.3. * 例。--- 设 \(\mathcal{H}\) 是维数为 \(n\) 的有限维实仿射空间 E 中的一个局部有限超平面集，C 是相对于 \(\mathcal{H}\) 的 E 的一个室（LIE, V, § 1, no. 3）。C 的闭包 \(\overline{\mathbf{C}}\) 是流形 E 的局部多面体子集；\(\overline{\mathbf{C}}\) 的正则边界是 C 的各个面的并（同上，no. 4）。特别地，若 C 是一个开单纯形（同上，no. 6），顶点为 \(a_0, \ldots, a_n\)，则 \(\overline{\mathbf{C}}\) 的边界是各个 \(\overline{\mathbf{C}_{\{i\}}}\)（\(0 \leqslant i \leqslant n\)）的并，其中 \(\mathbf{C}_{\{i\}}\) 是由顶点 \(a_0, \ldots, a_{i-1}, a_{i+1}, \ldots, a_n\) 在这些顶点生成的仿射空间中的开单纯形。*
11.3.4. 设 A 是 X 的局部多面体子集，\(\omega\) 是 X 上取值于 Banach 空间 E 的 \(n - 1\) 次扭曲微分形式；假设 \(\omega\) 属于类 \(\mathbf{C}^1\)，且其支集与 A 的交是紧的。于是，A（相应地 \(A'\)、\(\mathring{A}\)）的示性函数关于 X 上由 \(d\omega\) 定义的向量测度都是本质可积的，并且有

\[
\int_ {\mathrm{A}} d \omega = \int_ {\mathrm{A} ^ {\prime}} d \omega = \int_ {\mathring{\mathrm{A}}} d \omega.
\]

形式 \(\omega|\partial\mathbf{A}\)（对于 \(\partial\mathbf{A}\) 的典范嵌入所带的典范定向，该嵌入被看作片段 \(\mathbf{A}'\) 的边界并嵌入 \(\mathbf{X}\)）在 \(\partial\mathbf{A}\) 上可积，并且有

\[
\int_ {\mathbf {A}} d \omega = \int_ {\partial \mathbf {A}} \omega \quad (\text{Stokes公式}).
\]

当 X 配备一个定向、\(\partial\mathbf{A}\) 配备相应的定向（11.2.1）时，如果 \(\omega\) 是 X 上取值于 E 的 \(n-1\) 次普通微分形式、属于类 \(\mathbf{C}^{1}\)，且 \(\mathbf{A}\cap\operatorname{Supp}\omega\) 是紧的，则此公式仍然成立。

